\documentclass[11pt,letterpaper]{article}
\usepackage[utf8]{inputenc}
\usepackage[spanish]{babel}
\usepackage[margin=2.5cm]{geometry}
\usepackage{graphicx}
\usepackage{float}
\usepackage{longtable}
\usepackage{booktabs}
\usepackage{amssymb}
\usepackage[table]{xcolor}
\usepackage{hyperref}
\usepackage{enumitem}
\usepackage{fancyhdr}
\usepackage{titlesec}
\usepackage{lastpage}

\graphicspath{{figuras/}{../03_Modelado/01_Contexto/}{../03_Modelado/02_iStar_SD/}{../03_Modelado/03_iStar_SR/}{../03_Modelado/04_Casos_Uso/}{../03_Modelado/05_Clases/}{../03_Modelado/06_Secuencia/}{../03_Modelado/07_Actividad/}{../03_Modelado/08_Estados/}{../03_Modelado/09_DFD/}{../03_Modelado/10_Componentes/}{../03_Modelado/11_Despliegue/}{../03_Modelado/12_Prototipos_Interfaz/}}

\definecolor{sigagreen}{HTML}{1D6B36}
\definecolor{rotuloblue}{HTML}{0070C0}

\addto\captionsspanish{\renewcommand{\tablename}{Tabla}}
\addto\captionsspanish{\renewcommand{\listtablename}{\'Indice de tablas}}
\addto\captionsspanish{\renewcommand{\listfigurename}{\'Indice de figuras}}

\setcounter{secnumdepth}{0}
\setcounter{tocdepth}{3}
\titleformat{\section}{\normalfont\bfseries\fontsize{14}{16}\selectfont}{}{0em}{}
\titleformat{\subsection}{\normalfont\bfseries\fontsize{12}{14}\selectfont}{}{0em}{}
\titleformat{\subsubsection}{\normalfont\bfseries\fontsize{12}{14}\selectfont}{}{0em}{}
\linespread{1.3}

\pagestyle{fancy}
\fancyhf{}
\fancyfoot[C]{P\'agina \thepage\ de \pageref{LastPage}}

\title{}
\date{}

\begin{document}

\begin{titlepage}
\centering
{\bfseries\Large UNIVERSIDAD T\'ECNICA ESTATAL DE QUEVEDO\par}
\vspace{2mm}
{\bfseries\large FACULTAD DE CIENCIAS DE LA COMPUTACI\'ON\par}
\vspace{2mm}
{\bfseries\large CARRERA DE SOFTWARE\par}
\vspace{10mm}
{\bfseries TEMA:\par}
\vspace{2mm}
Proyecto Fin de Curso: Entrega Final (2B) -- Especificaci\'on de Requisitos\\
de Software (ERS/SRS) Consolidada + Componente Emp\'irico Ejecutado.
\vspace{8mm}

{\bfseries ESTUDIANTES:\par}
\vspace{2mm}
Mu\~noz Qui\~nonez Yeranick Esther \textbar\ CI: 1207929645 \textbar\ ymunozq@uteq.edu.ec\\
S\'anchez Cornejo Gary Alberto \textbar\ CI: 1208338291 \textbar\ gsanchezc6@uteq.edu.ec
\vspace{8mm}

{\bfseries CURSO:\par}
4TO "B"
\vspace{6mm}

{\bfseries DOCENTE:\par}
Ing. Guerrero Ulloa Gleiston Ciceron, PhD
\vspace{6mm}

{\bfseries MATERIA:\par}
Ingenier\'ia de Requerimientos
\vspace{6mm}

{\bfseries REPOSITORIO GITHUB:\par}
\url{https://github.com/gsanchezc6-beep/SIGA_FGMMN_ISR401_AVANCE_2B}
\vspace{6mm}

{\bfseries VERSI\'ON DE ESTE DOCUMENTO:\par}
\vspace{2mm}
Versi\'on \textbf{4.3} \textbar\ 31/08/2026 \textbar\ commit \texttt{0bc376d}\\
L\'inea base etiquetada \texttt{2B-final} en el repositorio declarado arriba.
\vspace{8mm}

{\bfseries QUEVEDO -- LOS R\'IOS -- ECUADOR\par}
{\bfseries 2026 -- 2027 PPA\par}
\end{titlepage}

\newpage

\section*{Resumen}
\addcontentsline{toc}{section}{Resumen}

Este documento especifica los requisitos del \textbf{Sistema Inteligente de Gestion de
Aulas (SIGA)}, una plataforma de Internet de las Cosas y aprendizaje automatico para la
Facultad de Ciencias de la Computacion de la Universidad Tecnica Estatal de Quevedo. La
especificacion se redacta conforme a la norma ISO/IEC/IEEE 29148:2018, y sus requisitos no
funcionales se cuantifican sobre el modelo de calidad de la norma ISO/IEC 25010:2023.

El levantamiento se apoya en \textbf{diez entrevistas semiestructuradas} con personal de
servicios generales, coordinacion de carrera y docencia, registradas en video y audio,
transcritas y codificadas tematicamente, y en un cuestionario aplicado a \textbf{sesenta}
respondientes. De ese material se derivan \textbf{veinticinco requisitos funcionales} y
\textbf{dieciseis no funcionales}, todos con umbral, unidad y metodo de comprobacion, mas
dieciseis casos de uso, diecisiete historias de usuario con criterios de aceptacion en
formato Dado-Cuando-Entonces, y dieciocho restricciones de diseno.

El sistema monitorea temperatura, humedad y ocupacion de las aulas, controla de forma
remota proyectores y climatizacion, genera alertas por anomalia, predice fallos de
equipamiento mediante aprendizaje automatico, gestiona solicitudes de mantenimiento y
produce reportes administrativos. Su \textbf{unico componente de inteligencia artificial}
es el analisis predictivo de fallos, especificado con requisitos de rendimiento, equidad,
explicabilidad y monitoreo en operacion, y clasificado como de riesgo minimo.

La especificacion se somete a una \textbf{auditoria de calidad} sobre seis metricas
--- completitud, consistencia, verificabilidad, trazabilidad, modificabilidad y correccion
--- cuyos conteos base se publican antes de calcular cada valor. La trazabilidad se
sostiene en una matriz que enlaza cada requisito con la evidencia de campo que lo origina,
sus casos de uso, historias, criterios de aceptacion y componentes de diseno.

\medskip
\noindent\textbf{Palabras clave.} Ingenieria de requisitos \textbullet{} Especificacion
de requisitos de software \textbullet{} ISO/IEC/IEEE 29148 \textbullet{} Internet de las
Cosas \textbullet{} Aulas inteligentes \textbullet{} Trazabilidad \textbullet{}
Requisitos de inteligencia artificial

\newpage

\section*{Abstract}
\addcontentsline{toc}{section}{Abstract}

This document specifies the requirements of the \textbf{Intelligent Classroom Management
System (SIGA)}, an Internet of Things and machine learning platform for the Faculty of
Computer Science at Universidad Tecnica Estatal de Quevedo, Ecuador. The specification
follows ISO/IEC/IEEE 29148:2018, and its non-functional requirements are quantified
against the ISO/IEC 25010:2023 product quality model.

Requirements elicitation draws on \textbf{ten semi-structured interviews} with facilities
staff, programme coordinators and teaching staff --- recorded on video and audio,
transcribed and thematically coded --- and on a questionnaire answered by \textbf{sixty}
respondents. From that material the document derives \textbf{twenty-five functional} and
\textbf{sixteen non-functional requirements}, each with a threshold, a unit and a
verification method, together with sixteen use cases, seventeen user stories with
Given-When-Then acceptance criteria, and eighteen design constraints.

The system monitors classroom temperature, humidity and occupancy, remotely controls
projectors and air conditioning, raises anomaly alerts, predicts equipment failures using
machine learning, manages maintenance requests and produces administrative reports. Its
\textbf{single artificial intelligence component} is the predictive failure analysis,
specified with performance, fairness, explainability and operational monitoring
requirements, and classified as minimal risk.

The specification undergoes a \textbf{quality audit} across six metrics --- completeness,
consistency, verifiability, traceability, modifiability and correctness --- whose base
counts are published before any value is computed. Traceability rests on a matrix linking
each requirement to the field evidence it originates from, together with its use cases,
user stories, acceptance criteria and design components.

\medskip
\noindent\textbf{Keywords.} Requirements engineering \textbullet{} Software requirements
specification \textbullet{} ISO/IEC/IEEE 29148 \textbullet{} Internet of Things
\textbullet{} Smart classrooms \textbullet{} Traceability \textbullet{} AI requirements

\newpage
\tableofcontents
\newpage
\listoftables
\newpage
\listoffigures
\newpage

\input{secciones_generadas.tex}

\end{document}
